\section{碰撞}
\begin{definition}[碰撞]
	两个或多个物体相互接近或发生接触时，在极短时间内，物体的运动状态发生显著变化。\
\end{definition}
除了球的撞击、打桩等以外，分子、原子等微观粒子之间的相互作用也是碰撞过程，被称为“散射”

\begin{definition}[正碰]
	对于两球，如果在碰撞前的速度在两球的中心线上，那么碰撞后的速度也都在这一连线上\
\end{definition}
在正碰过程中，由动量守恒定律可以得到$$
m_1v_{10}+m_2v_{20}=m_1v_1+m_2v_2   (1)
$$
\subsection{恢复系数}
碰撞定律：碰撞后两球的分离速度（$v_2-v_1$）与碰撞前的接近速度（$v_{10}-v_{20}$）成正比,
而这一比值由两球的材料性质决定，即：
$$
e=\dfrac{v_2-v_1}{v_{10}-v_{20}}   (2)
$$
通常把e叫做恢复系数
依据不同的e的取值，我们把正碰分为三种情况：
\begin{itemize}
    \item $e=0$:完全非弹性碰撞
    \item $0<e<1$:非弹性碰撞
    \item $e=1$:完全弹性碰撞
    
\end{itemize}
由（1）和（2）式联立可得：
$$
v_1=v_{10}-\dfrac{(1+e)m_2(v_{10}-v_{20})}{m_1+m_2}
$$
$$
v_2=v_{20}+\dfrac{(1+e)m_1(v_{10}-v_{20})}{m_1+m_2}
$$
\subsection{完全弹性碰撞}
令e=1可以得到
$$
v_1=\dfrac{(m_1-m_2)v_{10}+2m_2v_{20}}{m_1+m_2}
$$
$$
v_2=\dfrac{(m_2-m_1)v_{20}+2m_1v_{10}}{m_1+m_2}
$$
\begin{itemize}
    \item 设两球质量相等,即$m_1=m_2$,代入上式可得$v_1=v_{20} v_2=v_{10}$,此时两球的速度和能量发生了转换.

    \item 设$v_{20}=0$可得$v_1=\dfrac{(m_1-m_2)v_{10}}{m_1+m_2} , v_2=\dfrac{2m_1v_{10}}{m_1+m_2}$
    \item 假设$m_2>>m_1$,此时有$v_1=-v_{10},v_2=0$
\end{itemize}
\subsection{完全非弹性碰撞}
由完全非弹性碰撞e=0可得
$$
v_1=v_2=\dfrac{m_1v_{10}+m_2v_{20}}{m_1+m_2}
$$
此时碰撞中损失的机械能到达最大值，为：
$$
\Delta E=\dfrac{1}{2}(1-e^2)\dfrac{m_1m_2}{m_1+m_2}(v_{10}-v_{20})^2
$$
若此时两个物体中有一个物体是静止的（$v_{20}=0$）,则有：
$$
\Delta E=\dfrac{1}{2}(1-e^2)\dfrac{m_1m_2}{m_1+m_2}v_{10}^2
$$